% ===========================================================================
%  E-Trike Vero Board Wiring — RT & SYS ESP32-S3
%  Generated: 2026-07-23
%  Compile:  pdflatex -interaction=nonstopmode vero-board-wiring.tex
%  Source:   docs/hardware/wiring-vero-board.md
% ===========================================================================

\documentclass[10pt,a4paper,oneside]{report}

% ── Packages ──────────────────────────────────────────────────────────
\usepackage[utf8]{inputenc}
\usepackage[T1]{fontenc}
\usepackage{lmodern}
\usepackage[inner=20mm,outer=15mm,top=18mm,bottom=15mm,footskip=10mm,headsep=6mm]{geometry}
\usepackage{amsmath,amssymb}
\usepackage{booktabs}
\usepackage{longtable}
\usepackage{array}

% ── Table spacing and behavior ──────────────────────────────────────
\setlength{\LTpre}{4pt plus 2pt minus 2pt}
\setlength{\LTpost}{4pt plus 2pt minus 2pt}
\setlength{\tabcolsep}{3pt}
\renewcommand{\arraystretch}{1.1}
\raggedbottom
\usepackage{listings}
\usepackage[table]{xcolor}
\usepackage{soul}
\usepackage[normalem]{ulem}
\usepackage{hyperref}
\usepackage{fancyhdr}
\usepackage{titlesec}
\usepackage{enumitem}
\usepackage{alltt}
\usepackage{makecell}
\usepackage{caption}
\usepackage{float}
\usepackage{parskip}
\usepackage{textcomp}
\usepackage{calc}

% ── Prevent overfull hboxes in narrow table columns ─────────────────
\setlength{\emergencystretch}{3em}
\tolerance=3000
\hbadness=10000

% ── Prevent orphaned titles and headers ─────────────────────────────
\widowpenalty=10000
\clubpenalty=10000
\displaywidowpenalty=10000
\brokenpenalty=10000
\predisplaypenalty=10000
\postdisplaypenalty=10000
\interlinepenalty=10000

% ── Unicode character support ─────────────────────────────────────────
\DeclareUnicodeCharacter{03A9}{\ensuremath{\Omega}}
\DeclareUnicodeCharacter{03C0}{\ensuremath{\pi}}
\DeclareUnicodeCharacter{03B4}{\ensuremath{\delta}}
\DeclareUnicodeCharacter{03C9}{\ensuremath{\omega}}
\DeclareUnicodeCharacter{03B1}{\ensuremath{\alpha}}
\DeclareUnicodeCharacter{03B2}{\ensuremath{\beta}}
\DeclareUnicodeCharacter{03B8}{\ensuremath{\theta}}
\DeclareUnicodeCharacter{2264}{\ensuremath{\leq}}
\DeclareUnicodeCharacter{2265}{\ensuremath{\geq}}
\DeclareUnicodeCharacter{00B0}{\ensuremath{{}^\circ}}
\DeclareUnicodeCharacter{00B1}{\ensuremath{\pm}}
\DeclareUnicodeCharacter{00D7}{\ensuremath{\times}}
\DeclareUnicodeCharacter{00B5}{\ensuremath{\mu}}
\DeclareUnicodeCharacter{2192}{\ensuremath{\rightarrow}}
\DeclareUnicodeCharacter{2190}{\ensuremath{\leftarrow}}
\DeclareUnicodeCharacter{2194}{\ensuremath{\leftrightarrow}}
\DeclareUnicodeCharacter{25CB}{\ensuremath{\circ}}
\DeclareUnicodeCharacter{2713}{$\checkmark$}

% ── Title formatting (11pt chapter/section headings, 10pt body) ───────
\titleformat{\chapter}[display]
  {\normalfont\bfseries\fontsize{11}{13}\selectfont}
  {}{0pt}{}
\titlespacing*{\chapter}{0pt}{-2em}{8pt}
\titleformat{\section}
  {\normalfont\bfseries\fontsize{11}{13}\selectfont}
  {}{0pt}{}
\titleformat{\subsection}
  {\normalfont\bfseries\fontsize{10}{12}\selectfont}
  {}{0pt}{}
\titleformat{\subsubsection}
  {\normalfont\bfseries\fontsize{10}{12}\selectfont}
  {}{0pt}{}

% ── Hyperref setup ────────────────────────────────────────────────────
\hypersetup{
  colorlinks=true,
  linkcolor=black,
  urlcolor=blue!60!black,
  citecolor=black,
  pdfauthor={E-Trike Project},
  pdftitle={E-Trike Vero Board Wiring — RT & SYS ESP32-S3},
  pdfsubject={Vero board wiring for RT and SYS ESP32-S3 controllers},
  pdfkeywords={CAN bus, embedded systems, ESP32-S3, GPIO, pinout, vero board, wiring},
  bookmarksnumbered=true,
  bookmarksopen=true,
  bookmarksopenlevel=1
}

% ── Code listing style with strikethrough support ──────────────────────
\lstdefinestyle{etrike}{
  basicstyle=\ttfamily\footnotesize,
  breaklines=true,
  breakatwhitespace=true,
  postbreak=\mbox{$\hookrightarrow$\space},
  frame=single,
  framesep=3pt,
  xleftmargin=4pt,
  xrightmargin=4pt,
  aboveskip=6pt,
  belowskip=6pt,
  showstringspaces=false,
  columns=flexible,
  keepspaces=true,
  tabsize=2,
  captionpos=b,
  escapechar=|
}
\lstset{style=etrike}

% ── Header / footer ───────────────────────────────────────────────────
\pagestyle{fancy}
\fancyhf{}
\fancyhead[L]{\footnotesize\textit{E-Trike Vero Board Wiring}}
\fancyhead[R]{\footnotesize\textit{v0.8.0-alpha}}
\fancyfoot[C]{\thepage}
\renewcommand{\headrulewidth}{0.4pt}
\renewcommand{\footrulewidth}{0pt}
\setlength{\headheight}{14pt}
\setlength{\headsep}{6mm}

% ── Convenience commands ──────────────────────────────────────────────
\newcommand{\canid}[1]{\texttt{0x#1}}
\newcommand{\canname}[1]{\textsc{#1}}
\newcommand{\gpio}[1]{\texttt{GPIO#1}}
\newcommand{\voltage}[1]{\ensuremath{#1\,\text{V}}}
\newcommand{\bitrate}[1]{\ensuremath{#1\,\text{kbit/s}}}
\newcommand{\frequency}[1]{\ensuremath{#1\,\text{Hz}}}
\newcommand{\mcu}[1]{\texttt{#1}}
\newcommand{\pinname}[1]{\texttt{#1}}
\newcommand{\iicaddr}[1]{\texttt{0x#1}}

\newcolumntype{P}[1]{>{\centering\arraybackslash}p{#1}}
\newcolumntype{L}[1]{>{\raggedright\arraybackslash}p{#1}}

\begin{document}

% ═══════════════════════════════════════════════════════════════════════
%  CHAPTER 1 — RT ESP32-S3 — CAN Gateway & Real-Time Controller
% ═══════════════════════════════════════════════════════════════════════

\chapter{RT ESP32-S3 — CAN Gateway \& Real-Time Controller}

\textbf{Role:} CAN gateway (high $\leftrightarrow$ low), real-time kinematics, steering state machine, PID speed control\\
\textbf{Board:} ESP32-S3-DevKitC-1 (N16R8 — 16 MB flash, 8 MB octal PSRAM)\\
\textbf{CAN modules:} 2 — TWAI (built-in, low bus) + MCP2515 (SPI, high bus)

% ──────────────────────────────────────────────────────────────────────
\section{RT — CAN Module Wiring}

\noindent\textit{Note: Both CAN transceiver modules (WCMCU-230 and MCP2515) are mounted directly on the vero board. All connections between the ESP32-S3 dev board and the transceiver modules are made via traces/wires on the vero board — not via external Dupont cables.}

\subsection{Module A: WCMCU-230 (Low Bus via TWAI)}

\begin{longtable}{|L{1cm}|L{1.5cm}|L{2.5cm}|L{2.8cm}|L{4.5cm}|}
\hline
\textbf{GPIO} & \textbf{J1 Pin} & \textbf{Signal} & \textbf{Connected To} & \textbf{Notes} \\
\hline
\endfirsthead
\hline
\textbf{GPIO} & \textbf{J1 Pin} & \textbf{Signal} & \textbf{Connected To} & \textbf{Notes} \\
\hline
\endhead
— & 1 (3V3) & 3.3\,V power & WCMCU-230 VCC & Vero board 3.3\,V rail \\
\hline
— & 22 (GND) & Ground & WCMCU-230 GND & Vero board ground plane \\
\hline
5 & 5 & CAN TX (TWAI) & WCMCU-230 CTX & Vero board trace \\
\hline
4 & 4 & CAN RX (TWAI) & WCMCU-230 CRX & Vero board trace \\
\hline
— & — & CAN\_H & Low bus backbone & 22 AWG, yellow \\
\hline
— & — & CAN\_L & Low bus backbone & 22 AWG, green \\
\hline
— & — & GND & Low bus backbone & 22 AWG, black \\
\hline
\end{longtable}

\begin{itemize}[nosep]
  \item 120\,$\Omega$ terminator: \textbf{ON} (RT is a bus endpoint on the low bus)
  \item Source: \texttt{rt::kCanLowTxGpio=5}, \texttt{rt::kCanLowRxGpio=4}, bitrate 500\,kbit/s
\end{itemize}

\subsection{Module B: MCP2515 SPI (High Bus)}

\begin{longtable}{|L{1cm}|L{2.2cm}|L{3cm}|L{5.5cm}|}
\hline
\textbf{GPIO} & \textbf{Signal} & \textbf{Connected To} & \textbf{Notes} \\
\hline
\endfirsthead
\hline
\textbf{GPIO} & \textbf{Signal} & \textbf{Connected To} & \textbf{Notes} \\
\hline
\endhead
15 & SPI SCK & MCP2515 SCK & Vero board trace \\
\hline
16 & SPI MOSI & MCP2515 SI & Vero board trace \\
\hline
17 & SPI MISO & MCP2515 SO & Vero board trace \\
\hline
18 & SPI CS & MCP2515 CS & Vero board trace \\
\hline
47 & INT & MCP2515 INT & Vero board trace \\
\hline
— & Module supply & MCP2515 VCC & Vero board 5\,V rail \\
\hline
— (J1-22, GND) & Ground & MCP2515 GND & Vero board ground plane \\
\hline
— & CAN\_H & High bus backbone & 22 AWG, orange \\
\hline
— & CAN\_L & High bus backbone & 22 AWG, blue \\
\hline
\end{longtable}

\begin{itemize}[nosep]
  \item No onboard termination — CANalyst-II Ch0 provides the high-bus endpoint termination during bench work
  \item Source: \texttt{rt::kSpiSckGpio=15}, \texttt{rt::kSpiMosiGpio=16}, \texttt{rt::kSpiMisoGpio=17}, \texttt{rt::kSpiCsGpio=18}, \texttt{rt::kMcpIntGpio=47}, bitrate 500\,kbit/s
\end{itemize}

% ──────────────────────────────────────────────────────────────────────
\section{RT — Inputs}

\begin{longtable}{|L{0.5cm}|L{2.8cm}|L{1cm}|L{2cm}|L{2.5cm}|L{1cm}|L{3.5cm}|}
\hline
\textbf{\#} & \textbf{Signal} & \textbf{GPIO} & \textbf{Type} & \textbf{Connected To} & \textbf{Status} & \textbf{Notes} \\
\hline
\endfirsthead
\hline
\textbf{\#} & \textbf{Signal} & \textbf{GPIO} & \textbf{Type} & \textbf{Connected To} & \textbf{Status} & \textbf{Notes} \\
\hline
\endhead
1 & CAN RX (low bus) & 4 & TWAI RX & SN65HVD230 CRX & \textbf{Active} & Low bus frames from SYS, MTR, EPS-C, SEB \\
\hline
2 & MCP2515 MISO & 17 & SPI & MCP2515 SO & \textbf{Active} & High bus RX data (Host commands) \\
\hline
3 & MCP2515 INT & 47 & Digital, falling & MCP2515 INT & \textbf{Active} & RX buffer ready interrupt \\
\hline
4 & Developer override button [latch] & 42 & Digital, active-low & Latch switch $\rightarrow$ GND & & Internal pull-up; J3 pin 6; conflicts with external JTAG MTMS \\
\hline
\end{longtable}

% ──────────────────────────────────────────────────────────────────────
\subsection{RT — Encoder Inputs}

\begin{longtable}{|L{0.5cm}|L{2.8cm}|L{1cm}|L{2cm}|L{2.5cm}|L{1cm}|L{3.5cm}|}
\hline
\textbf{\#} & \textbf{Signal} & \textbf{GPIO} & \textbf{Type} & \textbf{Connected To} & \textbf{Status} & \textbf{Notes} \\
\hline
\endfirsthead
\hline
\textbf{\#} & \textbf{Signal} & \textbf{GPIO} & \textbf{Type} & \textbf{Connected To} & \textbf{Status} & \textbf{Notes} \\
\hline
\endhead
1 & Rear motor encoder A & 1 & PCNT input & Encoder ChA via terminal block & \textbf{Active} & Terminal block $\rightarrow$ GPIO1 \\
\hline
2 & Rear motor encoder B & 2 & PCNT input & Encoder ChB via terminal block & \textbf{Active} & Terminal block $\rightarrow$ GPIO2 \\
\hline
3 & Front wheel encoder A & 10 & PCNT input & Encoder ChA via terminal block & \textbf{Active} & Terminal block $\rightarrow$ GPIO10 \\
\hline
4 & Front wheel encoder B & 6 & PCNT input & Encoder ChB via terminal block & \textbf{Active} & Terminal block $\rightarrow$ GPIO6 \\
\hline
5 & Rear-left encoder A & 9 & PCNT input & Encoder ChA via terminal block & \textbf{Active} & Terminal block $\rightarrow$ GPIO9 \\
\hline
6 & Rear-left encoder B & 12 & PCNT input & Encoder ChB via terminal block & \textbf{Active} & Terminal block $\rightarrow$ GPIO12 \\
\hline
7 & Rear-right encoder A & 13 & PCNT input & Encoder ChA via terminal block & \textbf{Active} & Terminal block $\rightarrow$ GPIO13 \\
\hline
8 & Rear-right encoder B & 14 & PCNT input & Encoder ChB via terminal block & \textbf{Active} & Terminal block $\rightarrow$ GPIO14 \\
\hline
\end{longtable}

\noindent\textit{Note: All encoder signals must first connect to a \textbf{Pluggable Screw Terminal Block Connector} on the vero board. Encoder wires (from harness) terminate in the screw terminal block; short Dupont jumper wires then run from the terminal block to the ESP32 GPIO pins.}

% ──────────────────────────────────────────────────────────────────────
\section{RT — Outputs}

\begin{longtable}{|L{0.5cm}|L{2.2cm}|L{1cm}|L{2cm}|L{2.5cm}|L{4.5cm}|}
\hline
\textbf{\#} & \textbf{Signal} & \textbf{GPIO} & \textbf{Type} & \textbf{Connected To} & \textbf{Notes} \\
\hline
\endfirsthead
\hline
\textbf{\#} & \textbf{Signal} & \textbf{GPIO} & \textbf{Type} & \textbf{Connected To} & \textbf{Notes} \\
\hline
\endhead
1 & CAN TX (low bus) & 5 & TWAI TX & SN65HVD230 CTX & 0x204, 0x205, 0x169, 0x001, 0x302, 0x7FD on low bus \\
\hline
2 & SPI SCK & 15 & SPI out & MCP2515 SCK & 10 MHz max \\
\hline
3 & SPI MOSI & 16 & SPI out & MCP2515 SI & High bus TX data \\
\hline
4 & SPI CS & 18 & Digital out & MCP2515 CS & Active low \\
\hline
5 & WDT Toggle & 21 & Digital out & TPS3850 WDI pin & Toggled at 100 Hz by control\_task \\
\hline
\end{longtable}

% ──────────────────────────────────────────────────────────────────────
\section{RT — GPIO Quick Reference}

\textbf{General GPIO:}
\begin{lstlisting}
GPIO 4  : CAN RX - low bus TWAI
GPIO 5  : CAN TX - low bus TWAI
GPIO 7  : (unused)
GPIO 8  : (unused)
GPIO 11 : (unused)
GPIO 15 : MCP2515 SCK
GPIO 16 : MCP2515 MOSI
GPIO 17 : MCP2515 MISO
GPIO 18 : MCP2515 CS
GPIO 21 : WDT toggle -> TPS3850 WDI
GPIO 42 : Developer override latch (active-low, J3 pin 6)
GPIO 47 : MCP2515 INT
\end{lstlisting}

\textbf{Encoder GPIO (via terminal block) — \st{not yet wired}:}
\begin{lstlisting}
GPIO 1  : |\sout{Rear motor encoder A (PCNT)}|
GPIO 2  : |\sout{Rear motor encoder B (PCNT)}|
GPIO 6  : |\sout{Front wheel encoder B (PCNT)}|
GPIO 9  : |\sout{Rear-left encoder A (PCNT)}|
GPIO 10 : |\sout{Front wheel encoder A (PCNT)}|
GPIO 12 : |\sout{Rear-left encoder B (PCNT)}|
GPIO 13 : |\sout{Rear-right encoder A (PCNT)}|
GPIO 14 : |\sout{Rear-right encoder B (PCNT)}|
\end{lstlisting}

% ──────────────────────────────────────────────────────────────────────
\section{RT — Power \& Wiring Notes}

\begin{longtable}{|L{1.5cm}|L{3.5cm}|L{6cm}|}
\hline
\textbf{Rail} & \textbf{Source} & \textbf{Used For} \\
\hline
\endfirsthead
\hline
\textbf{Rail} & \textbf{Source} & \textbf{Used For} \\
\hline
\endhead
3.3\,V & USB or dev board regulator & ESP32-S3, WCMCU-230 transceiver \\
\hline
5\,V & USB or dev board regulator & MCP2515 module VCC \\
\hline
GND & USB or dev board & Common ground — shared with both CAN modules \\
\hline
\end{longtable}

\textbf{Encoder Wiring:} All encoder signals (GPIO 1, 2, 6, 9, 10, 12, 13, 14) must route through \textbf{Pluggable Screw Terminal Block Connectors} mounted on the vero board.

\newpage

% ═══════════════════════════════════════════════════════════════════════
%  CHAPTER 2 — SYS ESP32-S3 — Safety & Body Controller
% ═══════════════════════════════════════════════════════════════════════

\chapter{SYS ESP32-S3 — Safety \& Body Controller}

\textbf{Role:} Safety monitor (EGAS Level 2), body control, and brake-by-wire interface.\\
Motor DAC and gear output are bench-only on SYS; no vehicle motor-actuation path is approved until MTR hardware is implemented and validated.\\
\textbf{Board:} ESP32-S3-DevKitC-1 (N16R8 — 16 MB flash, 8 MB octal PSRAM)\\
\textbf{CAN modules:} 1 — TWAI (built-in, low bus only)

% ──────────────────────────────────────────────────────────────────────
\section{SYS — CAN Module Wiring}

\noindent\textit{Note: The WCMCU-230 CAN transceiver module is mounted directly on the vero board. All connections between the ESP32-S3 dev board and the transceiver module are made via traces/wires on the vero board — not via external Dupont cables.}

\subsection{Module: WCMCU-230 (Low Bus via TWAI)}

\begin{longtable}{|L{1cm}|L{1.5cm}|L{2.5cm}|L{2.8cm}|L{4.5cm}|}
\hline
\textbf{GPIO} & \textbf{J1 Pin} & \textbf{Signal} & \textbf{Connected To} & \textbf{Notes} \\
\hline
\endfirsthead
\hline
\textbf{GPIO} & \textbf{J1 Pin} & \textbf{Signal} & \textbf{Connected To} & \textbf{Notes} \\
\hline
\endhead
— & 1 (3V3) & 3.3\,V power & WCMCU-230 VCC & Vero board 3.3\,V rail \\
\hline
— & 22 (GND) & Ground & WCMCU-230 GND & Vero board ground plane \\
\hline
5 & 5 & CAN TX (TWAI) & WCMCU-230 CTX & Vero board trace \\
\hline
4 & 4 & CAN RX (TWAI) & WCMCU-230 CRX & Vero board trace \\
\hline
— & — & CAN\_H & Low bus backbone & 22 AWG, yellow \\
\hline
— & — & CAN\_L & Low bus backbone & 22 AWG, green \\
\hline
— & — & GND & Low bus backbone & 22 AWG, black \\
\hline
\end{longtable}

\begin{itemize}[nosep]
  \item 120\,$\Omega$ terminator: \textbf{ON} (SYS is a bus endpoint on the low bus)
  \item Source: \texttt{sys::kCanTxGpio=5}, \texttt{sys::kCanRxGpio=4}, bitrate 500\,kbit/s
\end{itemize}

% ──────────────────────────────────────────────────────────────────────
\section{SYS — Safety Inputs}

\begin{longtable}{|L{0.5cm}|L{2.8cm}|L{1cm}|L{2cm}|L{2.5cm}|L{4.5cm}|}
\hline
\textbf{\#} & \textbf{Signal} & \textbf{GPIO} & \textbf{Type} & \textbf{Connected To} & \textbf{Notes} \\
\hline
\endfirsthead
\hline
\textbf{\#} & \textbf{Signal} & \textbf{GPIO} & \textbf{Type} & \textbf{Connected To} & \textbf{Notes} \\
\hline
\endhead
1 & ESTOP Button & 1 & Digital, NC, active-low & NC contact from 3.3\,V to GPIO1 & 10k external pull-down at GPIO1. LOW = ESTOP, including an open wire. \\
\hline
2 & Brake Lever & 2 & Digital, NO, active-low & Lever switch $\rightarrow$ GND & Internal pull-up. LOW = pressed $\rightarrow$ SEB via CAN 0x7B9 \\
\hline
3 & START Button & 41 & Digital, NO, active-low & Green momentary $\rightarrow$ GND & Internal pull-up. Press exits ESTOP to MANUAL. \\
\hline
4 & MODE Button & 11 & Digital, NO, active-low & Momentary $\rightarrow$ GND & Internal pull-up. Toggles MANUAL $\leftrightarrow$ AUTO. Long-press 3s exits ESTOP. \\
\hline
5 & Ignition relay & 8 & Reserved output & Not implemented by production firmware & Do not wire as a vehicle power-control path. \\
\hline
6 & Developer override button [latch] & 42 & Digital, active-low & Latch switch $\rightarrow$ GND & Internal pull-up; J3 pin 6; conflicts with external JTAG MTMS. \\
\hline
7 & CAN RX (low bus) & 4 & TWAI RX & SN65HVD230 CRX & RT commands, MTR feedback (0x206), SEB feedback (0x721) \\
\hline
\end{longtable}

% ──────────────────────────────────────────────────────────────────────
\section{SYS — Signal Lights \& Switches}

\begin{longtable}{|L{0.5cm}|L{2.8cm}|L{1cm}|L{2cm}|L{2.5cm}|L{4.5cm}|}
\hline
\textbf{\#} & \textbf{Signal} & \textbf{GPIO} & \textbf{Type} & \textbf{Connected To} & \textbf{Notes} \\
\hline
\endfirsthead
\hline
\textbf{\#} & \textbf{Signal} & \textbf{GPIO} & \textbf{Type} & \textbf{Connected To} & \textbf{Notes} \\
\hline
\endhead
1 & Left Turn Switch & 9 & Digital, NO, active-low & Handlebar momentary $\rightarrow$ GND & Internal pull-up \\
\hline
2 & Right Turn Switch & 6 & Digital, NO, active-low & Handlebar momentary $\rightarrow$ GND & Internal pull-up \\
\hline
3 & Headlight Switch & 7 & Digital, NO, active-low & Handlebar toggle $\rightarrow$ GND & Internal pull-up \\
\hline
4 & Left Turn Lamp & 18 & Digital out & Relay coil $\rightarrow$ 12\,V & Relay COM = 12\,V, NO = lamp \\
\hline
5 & Right Turn Lamp & 19 & Digital out & Relay coil $\rightarrow$ 12\,V & Relay COM = 12\,V, NO = lamp \\
\hline
6 & Brake Light & 21 & Digital out & Relay coil $\rightarrow$ 12\,V & Relay COM = 12\,V, NO = lamp \\
\hline
7 & Headlight & 10 & Digital out & Relay coil driver $\rightarrow$ 12\,V & ESP GPIO must not drive the relay coil directly. \\
\hline
\end{longtable}

% ──────────────────────────────────────────────────────────────────────
\section{SYS — Indicators \& Power Control}

\begin{longtable}{|L{0.5cm}|L{2.8cm}|L{1cm}|L{2cm}|L{2.5cm}|L{4.5cm}|}
\hline
\textbf{\#} & \textbf{Signal} & \textbf{GPIO} & \textbf{Type} & \textbf{Connected To} & \textbf{Notes} \\
\hline
\endfirsthead
\hline
\textbf{\#} & \textbf{Signal} & \textbf{GPIO} & \textbf{Type} & \textbf{Connected To} & \textbf{Notes} \\
\hline
\endhead
1 & AUTO Mode Bulb & 48 & Digital out & Lamp driver $\rightarrow$ 12\,V & ESP GPIO must not drive the lamp directly. \\
\hline
2 & MANUAL Mode Bulb & 39 & Digital out & Lamp driver $\rightarrow$ 12\,V & Conflicts with external JTAG use. \\
\hline
3 & READY Bulb & 17 & Digital out & Relay coil $\rightarrow$ 12\,V & Green — system ready (AUTO/MANUAL, RT alive, no faults) \\
\hline
4 & ESTOP Bulb & 20 & Digital out & Relay coil $\rightarrow$ 12\,V & Red — dedicated ESTOP indicator \\
\hline
5 & Bypass Indicator & 14 & Digital out & Lamp driver / LED $\rightarrow$ GND & Yellow/amber — developer override / safety bypass active \\
\hline
6 & 12V Power Relay & 40 & Digital out & Relay driver $\rightarrow$ 12\,V & Accessory relay, opens in ESTOP. Conflicts with external JTAG use. \\
\hline
7 & WDT Toggle & 47 & Digital out & TPS3850 WDI pin & Toggled at 20 Hz by safety task. \\
\hline
\end{longtable}

% ──────────────────────────────────────────────────────────────────────
\section{SYS — GPIO Quick Reference}

\begin{lstlisting}
GPIO 1  : ESTOP button (NC from 3.3 V, external pull-down, active-low)
GPIO 2  : Brake lever (active-low) -> SEB via CAN 0x7B9
GPIO 4  : CAN RX - low bus TWAI
GPIO 5  : CAN TX - low bus TWAI
GPIO 6  : |\sout{Right turn switch (active-low)}|
GPIO 7  : |\sout{Headlight switch (active-low)}|
GPIO 8  : Reserved ignition output - not implemented
GPIO 9  : |\sout{Left turn switch (active-low)}|
GPIO 10 : |\sout{Headlight relay driver}|
GPIO 11 : MODE button (active-low) - publishes CAN 0x110
GPIO 12-13 : Bench-only gear sense
GPIO 14 : Bypass indicator (yellow/amber)
GPIO 15-16 : Bench-only MCP4725 I2C
GPIO 17 : |\sout{READY bulb (green)}|
GPIO 18 : |\sout{Left turn lamp}|
GPIO 19 : |\sout{Right turn lamp}|
GPIO 20 : |\sout{ESTOP bulb (red)}|
GPIO 21 : |\sout{Brake light}|
GPIO 33-37 : Unavailable on N16R8 octal-memory modules
GPIO 39 : MANUAL mode bulb - external JTAG conflict
GPIO 40 : 12V accessory relay - external JTAG conflict
GPIO 41 : START button (active-low) - external JTAG conflict
GPIO 42 : Developer override latch (active-low, J3 pin 6) - external JTAG conflict
GPIO 47 : WDT toggle -> TPS3850 WDI (20 Hz)
GPIO 48 : AUTO mode bulb
\end{lstlisting}

% ──────────────────────────────────────────────────────────────────────
\section{SYS — Power}

\begin{longtable}{|L{1.5cm}|L{3.5cm}|L{6cm}|}
\hline
\textbf{Rail} & \textbf{Source} & \textbf{Used For} \\
\hline
\endfirsthead
\hline
\textbf{Rail} & \textbf{Source} & \textbf{Used For} \\
\hline
\endhead
3.3\,V & USB or dev board regulator & ESP32-S3, WCMCU-230 transceiver \\
\hline
5\,V & USB or dev board regulator & MCP4725 DAC VCC (for 0–5\,V output) \\
\hline
12\,V & DC-DC converter (via 12V relay) & Lamps, bulbs, relay coils \\
\hline
GND & USB or dev board & Common ground \\
\hline
\end{longtable}

\end{document}
